%% kapa: exam=1 round=29 year=2018 period=2 form=A subject=appraisal_law subject_ko=감정평가관계법규 seq=28 number=28
%% source: Q-Net 감정평가사 기출문제 (cst00309) artlSeq=5202808 file=extracted/1008/5202808/2018년 제29회 감정평가사 1차시험 기출문제/2018년 제29회 감정평가사 1차시험 기출문제/2교시 A형.pdf
%% conversion: kordoc 4.24.51 -> tools/gen1.py (auto); review=proofread-sonnet; flags=none
\begin{question}{28}
공간정보의 구축 및 관리 등에 관한 법령상 지번의 부여 등에 관한 설명으로 옳지 않은 것은?
\begin{choices}
  \item 지번은 지적소관청이 지번부여지역별로 차례대로 부여한다.
  \item 지번은 북서에서 남동으로 순차적으로 부여한다.
  \item 지번변경 승인신청을 받은 승인권자는 지번변경 사유 등을 심사한 후 그 결과를 지적소관청에 통지하여야 한다.
  \item 지번은 아라비아숫자로 표기하되, 임야대장 및 임야도에 등록하는 토지의 지번은 숫자 앞에 "산"자를 붙인다.
  \item 지적소관청이 지적공부에 등록된 지번을 변경하려면 국토교통부장관의 승인을 받아야 한다.
\end{choices}
\end{question}
\answer{5}
